% Source: Gerzensee "Real Analysis" slides 12-15 (Matrix Representation Theorem,
% Jacobian matrix, determinant as area magnification factor).
% Supplementary: Fukuda (2026), Ch. 4 (Basic Operations of Matrices, pp. 135-156),
% Ch. 5 (Fundamentals of Vector Spaces, pp. 157-184), and Sec. 10.7 (Quadratic
% Forms, pp. 363-365).
% External references actually cited in this chapter: Axler (2024),
% Ljungqvist-Sargent (2018).

\chapter{Vector Spaces, Linear Maps, and Matrices}\label{ch:linear}

\section{Motivation: Linearity as the First Approximation}\label{sec:lin-motivation}

Linear algebra enters economics along two distinct routes, and it is worth
separating them because they require different parts of the theory.

The first route is exact. Some economic models are linear by construction: a
Leontief input--output system, a linear-quadratic control problem, a system of
linear difference equations describing the dynamics of a linearised
macroeconomic model, the normal equations of ordinary least squares. Here one
needs solvability of $Ax=b$, invertibility, rank, and --- for dynamics --- the
spectrum of $A$.

The second route is approximate, and it is the more pervasive. The central idea
of differential calculus is that a smooth function is, near a point, well
approximated by a linear one; the derivative \emph{is} that linear map. The
Gerzensee slides make this explicit when they present the total differential
$\diff y=\ip{\nabla f}{\diff x}$ as an inner product and the Jacobian as the matrix
representing $df$ \autocite[slides 10--15]{fukuda2026analysis}. Every
comparative-statics computation in economics is the solution of a linear system
obtained in this way, and the invertibility condition in the implicit function
theorem (\autoref{thm:implicit-function}) is exactly the requirement that the
relevant linear system be uniquely solvable.

This chapter develops the theory in the coordinate-free language of vector spaces
and linear maps, with matrices appearing as coordinate representations. The
payoff is that the same theorems apply verbatim in the infinite-dimensional
function spaces of \autoref{ch:completeness} and \autoref{ch:dynamic}, where no
matrix exists.

\section{Vector Spaces}\label{sec:lin-vs}

\begin{definition}[Vector space]\label{def:vector-space}
	A \dfn{vector space} over $\R$ is a set $V$ with an addition
	$+\colon V\times V\to V$ and a scalar multiplication
	$\cdot\colon\R\times V\to V$ such that $(V,+)$ is an abelian group with neutral
	element $0$, and for all $\lambda,\mu\in\R$ and $v,w\in V$,
	\[
		\lambda(v+w)=\lambda v+\lambda w,\quad
		(\lambda+\mu)v=\lambda v+\mu v,\quad
		(\lambda\mu)v=\lambda(\mu v),\quad
		1\cdot v=v.
	\]
	Elements of $V$ are \dfnas{vector}{vector}s.
\end{definition}

\begin{eg}[Vector spaces used in this volume]\label{eg:vector-spaces}
	\begin{enumerate}[(i)]
		\item $\R^n$ with coordinatewise operations. This is the space of commodity
		      bundles, price vectors, and parameter vectors.
		\item The set $M_{m\times n}(\R)$ of $m\times n$ real matrices.
		\item The set $\R^S$ of all functions $f\colon S\to\R$ on a set $S$, with
		      pointwise operations. When $S=\{1,\dots,n\}$ this is $\R^n$; when
		      $S=\N$ it is the space of real sequences, the natural home of an
		      infinite consumption path $(c_t)_{t\in\N_0}$; when $S=\Omega$ is a sample
		      space it is the space of random variables.
		\item The set $\Cspace{[a,b]}{\R}$ of continuous real functions on $[a,b]$, and
		      the set $\Bspace{X}{\R}$ of bounded real functions on a set $X$. These are
		      the spaces on which the Bellman operator of \autoref{ch:dynamic} acts.
	\end{enumerate}
	Items (iii) and (iv) are infinite-dimensional. That is precisely why the
	coordinate-free treatment is worth the effort.
\end{eg}

\begin{definition}[Subspace, span, linear independence, basis]\label{def:span-basis}
	Let $V$ be a vector space over $\R$.
	\begin{enumerate}[(i)]
		\item A \dfn{subspace} of $V$ is a non-empty $W\subseteq V$ closed under
		      addition and scalar multiplication. Every subspace contains $0$.
		\item The \dfn{span} of $S\subseteq V$ is
		      $\Span S\eqdef\bigl\{\sum_{i=1}^k\lambda_iv_i : k\in\N,\ \lambda_i\in\R,\
			      v_i\in S\bigr\}$, with $\Span\emptyset\eqdef\{0\}$. It is the smallest
		      subspace containing $S$.
		\item A finite list $v_1,\dots,v_k$ is \dfnas{linearly independent}{linear
			      independence} if $\sum_i\lambda_iv_i=0$ implies
		      $\lambda_1=\dots=\lambda_k=0$, and \dfnas{linearly dependent}{linear
			      dependence} otherwise.
		\item A \dfn{basis} of $V$ is a linearly independent list spanning $V$. If $V$
		      has a finite basis it is \dfnas{finite-dimensional}{finite-dimensional
			      space} and its \dfn{dimension} $\dim V$ is the common length of all its
		      bases.
	\end{enumerate}
\end{definition}

\begin{theorem}[Well-definedness of dimension]\label{thm:dimension-well-defined}
	In a finite-dimensional vector space, every two bases have the same length, and
	every linearly independent list can be extended to a basis.
\end{theorem}

\begin{proof}
	The engine is the exchange lemma: in a vector space spanned by a list of length
	$m$, every linearly independent list has length at most $m$. Given that, if
	$B_1$ and $B_2$ are bases of lengths $m_1$ and $m_2$, then $B_1$ is independent
	and $B_2$ spans, giving $m_1\leq m_2$; exchanging roles gives $m_2\leq m_1$.

	For the exchange lemma itself, let $w_1,\dots,w_m$ span $V$ and let
	$v_1,\dots,v_k$ be independent with $k>m$. Adjoin $v_1$ to the spanning list:
	the list $v_1,w_1,\dots,w_m$ is dependent (since $v_1\in\Span\{w_j\}$), so some
	$w_j$ lies in the span of its predecessors and can be deleted, leaving a
	spanning list of length $m$ containing $v_1$. Repeating with
	$v_2,\dots,v_k$ --- at each step the deleted vector must be some $w_j$ rather
	than a previously inserted $v_i$, because the $v_i$ are independent --- exhausts
	the $w_j$ after $m$ steps while $v_{m+1}$ still remains, a contradiction. The
	extension statement follows by adjoining vectors outside the current span until
	the span is all of $V$, which terminates by the same length bound. A complete
	treatment is \textcite[Ch.~2]{axler2024linear}.
\end{proof}

\section{Linear Maps and their Matrix Representation}\label{sec:lin-maps}

\begin{definition}[Linear map]\label{def:linear-map}
	Let $V,W$ be vector spaces over $\R$. A map $T\colon V\to W$ is \dfn{linear} if
	$T(\lambda v+\mu v')=\lambda T(v)+\mu T(v')$ for all $\lambda,\mu\in\R$ and
	$v,v'\in V$. Its \dfn{kernel} and \dfn{image} are
	\[
		\Ker T\eqdef\{v\in V : T(v)=0\},\qquad
		\Image T\eqdef\{T(v) : v\in V\},
	\]
	subspaces of $V$ and $W$ respectively. The \dfn{rank} of $T$ is
	$\rank T\eqdef\dim\Image T$.
\end{definition}

\begin{theorem}[Matrix representation theorem]\label{thm:matrix-rep}
	A map $T\colon\R^n\to\R^m$ is linear if and only if there exists a unique
	$m\times n$ matrix $A$ with $T(x)=Ax$ for every $x\in\R^n$. Moreover
	\[
		A=\bigl[\,T(e^1)\ \ T(e^2)\ \cdots\ T(e^n)\,\bigr],
	\]
	where $e^i$ is the $i$-th standard basis vector, so that the $i$-th column of
	$A$ is $Ae^i=T(e^i)$.
\end{theorem}

\begin{proof}
	If $T(x)=Ax$ then linearity is the distributive law for matrix multiplication.

	Conversely, let $T$ be linear and define $A$ by the displayed formula. Writing
	$x=\sum_{i=1}^n x_ie^i$ and using linearity $n-1$ times,
	\[
		T(x)=T\Bigl(\sum_{i=1}^n x_ie^i\Bigr)=\sum_{i=1}^n x_i\,T(e^i)
		=\bigl[\,T(e^1)\ \cdots\ T(e^n)\,\bigr]\begin{pmatrix}x_1\\\vdots\\x_n\end{pmatrix}
		=Ax .
	\]
	For uniqueness, suppose $Ax=Bx$ for all $x$. Taking $x=e^i$ gives
	$Ae^i=Be^i$, that is, the $i$-th columns of $A$ and $B$ agree; as this holds for
	every $i$, $A=B$.
\end{proof}

This is Theorem~5.1 of \textcite{fukuda2026notes} and is presented on slides
12--14 of \autocite{fukuda2026analysis}. Two readings of $Ax$ are useful and
both appear later. Row by row, $(Ax)_i=\ip{a^i}{x}$ where $a^i$ is the $i$-th row
of $A$; this is the reading under which each row is a linear functional, used in
\autoref{ch:metric} for the Riesz representation theorem. Column by column,
$Ax=\sum_i x_i A e^i$ is a linear combination of the columns of $A$ with weights
$x_i$; this is the reading under which $\Image A$ is the column space, used in
\autoref{ch:projection} for least squares.

\begin{theorem}[Rank--nullity]\label{thm:rank-nullity}
	Let $V$ be finite-dimensional and $T\colon V\to W$ linear. Then
	\[
		\dim V=\dim\Ker T+\rank T .
	\]
\end{theorem}

\begin{proof}
	Let $u_1,\dots,u_k$ be a basis of $\Ker T$ and extend it, by
	\autoref{thm:dimension-well-defined}, to a basis
	$u_1,\dots,u_k,v_1,\dots,v_r$ of $V$, so $\dim V=k+r$. We claim
	$T(v_1),\dots,T(v_r)$ is a basis of $\Image T$.

	They span: any $T(v)$ with $v=\sum_i\alpha_iu_i+\sum_j\beta_jv_j$ equals
	$\sum_j\beta_jT(v_j)$ because $T(u_i)=0$. They are independent: if
	$\sum_j\beta_jT(v_j)=0$ then $T\bigl(\sum_j\beta_jv_j\bigr)=0$, so
	$\sum_j\beta_jv_j\in\Ker T$ and hence $\sum_j\beta_jv_j=\sum_i\alpha_iu_i$ for
	some $\alpha_i$; independence of the full basis forces all $\beta_j$ (and
	$\alpha_i$) to vanish. Therefore $\rank T=r$ and $\dim V=k+r$.
\end{proof}

\begin{corollary}[Solvability of $Ax=b$]\label{cor:solvability}
	Let $A$ be $m\times n$ and $b\in\R^m$.
	\begin{enumerate}[(i)]
		\item $Ax=b$ has a solution if and only if $b\in\Image A$, equivalently
		      $\rank[A\ b]=\rank A$.
		\item If a solution exists, the solution set is $x_0+\Ker A$ for any particular
		      solution $x_0$; it is a singleton if and only if $\rank A=n$.
		\item For $m=n$: $A$ is invertible $\iff$ $\rank A=n$ $\iff$ $\Ker A=\{0\}$
		      $\iff$ $\det A\neq0$ $\iff$ $Ax=b$ has exactly one solution for every $b$.
	\end{enumerate}
\end{corollary}

\begin{proof}
	(i) is the definition of the image. (ii) If $Ax_0=b$ then $Ax=b$ iff
	$A(x-x_0)=0$. The solution is unique iff $\Ker A=\{0\}$, which by
	\autoref{thm:rank-nullity} is $\rank A=n$. (iii) For $m=n$,
	\autoref{thm:rank-nullity} makes $\Ker A=\{0\}$ and $\Image A=\R^n$ equivalent;
	the equivalence with $\det A\neq0$ is the standard determinant criterion, and
	invertibility is then existence and uniqueness for every $b$.
\end{proof}

\begin{remark}[Square matters]\label{rem:square-matters}
	Part (iii) is a statement about square matrices only. For $m\neq n$ injectivity
	and surjectivity are genuinely different, and the distinction is economically
	substantive: in the least-squares problem of \autoref{ch:projection} the design
	matrix $X$ is $n\times k$ with $n\gg k$, so $Xb=y$ typically has no solution ---
	which is exactly why one projects instead of solving. In econometrics the
	condition $\rank X=k$ is the no-perfect-multicollinearity assumption, and by
	\autoref{cor:solvability}(ii) it is precisely what makes the estimator unique.
\end{remark}

\section{Determinants and Volume}\label{sec:lin-det}

\begin{definition}[Determinant]\label{def:determinant}
	The \dfn{determinant} is the unique function
	$\det\colon M_{n\times n}(\R)\to\R$ that is linear in each column separately,
	vanishes when two columns coincide, and satisfies $\det I_n=1$. Explicitly
	\[
		\det A=\sum_{\sigma\in S_n}\sgn(\sigma)\prod_{i=1}^n a_{\sigma(i)i},
	\]
	the sum running over permutations of $\{1,\dots,n\}$.
\end{definition}

\begin{proposition}[Properties]\label{prop:det-properties}
	For $A,B\in M_{n\times n}(\R)$: $\det(AB)=\det A\det B$;
	$\det A^{\transp}=\det A$; $A$ is invertible if and only if $\det A\neq0$, in
	which case $\det(A^{-1})=(\det A)^{-1}$; and for $\lambda\in\R$,
	$\det(\lambda A)=\lambda^n\det A$.
\end{proposition}

The geometric content is what the Gerzensee slides emphasise: $\abs{\det A}$ is
the factor by which the linear map $x\mapsto Ax$ magnifies $n$-dimensional
volume, and the sign of $\det A$ records whether orientation is preserved
\autocite[slide 15]{fukuda2026analysis}. Two consequences are used later.

\begin{eg}[Change of variables]\label{eg:det-jacobian}
	In the change-of-variables formula for multiple integrals
	(\autoref{thm:change-of-variables}), the factor $\abs{\det J}$ appears exactly
	because $J$ is the linear approximation of the substitution and $\abs{\det J}$
	is its volume magnification. The polar-coordinate computation of the Gaussian
	integral in \autocite[slides 91--93]{fukuda2026calculus}, where
	$\diff x\,\diff y=r\,\diff r\,\diff \theta$, is this statement with
	\[
		J=\begin{pmatrix}\partial x/\partial r & \partial x/\partial\theta\\
			\partial y/\partial r & \partial y/\partial\theta\end{pmatrix}
		=\begin{pmatrix}\cos\theta & -r\sin\theta\\ \sin\theta & r\cos\theta\end{pmatrix},
		\qquad \det J=r(\cos^2\theta+\sin^2\theta)=r>0 .
	\]
	The full derivation is given in \autoref{sec:int-gaussian}.
\end{eg}

\begin{eg}[Cramer's rule and comparative statics]\label{eg:cramer}
	If $\det A\neq0$, the unique solution of $Ax=b$ has $i$-th coordinate
	$x_i=\det A_i/\det A$, where $A_i$ is $A$ with its $i$-th column replaced by
	$b$. In comparative statics the system $Ax=b$ arises by differentiating
	equilibrium conditions, $x$ collects the responses of endogenous variables, and
	$b$ the direct effect of the parameter; Cramer's rule then delivers the sign of
	an individual response from the signs of two determinants. This is how sign
	restrictions on a Hessian or a Jacobian translate into comparative-statics
	predictions in \autoref{ch:calculusn}.
\end{eg}

\section{Eigenvalues, Diagonalisation, and Symmetric Matrices}\label{sec:lin-spectral}

\begin{definition}[Eigenvalue and eigenvector]\label{def:eigen}
	Let $A\in M_{n\times n}(\R)$. A scalar $\lambda\in\C$ is an \dfn{eigenvalue} of
	$A$ if $Av=\lambda v$ for some $v\neq0$; such a $v$ is an \dfn{eigenvector}. The
	eigenvalues are the roots of the \dfn{characteristic polynomial}
	$p_A(\lambda)=\det(A-\lambda I)$. The \dfn{spectral radius} is
	$\rho(A)\eqdef\max\{\abs{\lambda} : \lambda\text{ an eigenvalue of }A\}$.
\end{definition}

\begin{theorem}[Spectral theorem for symmetric matrices]\label{thm:spectral}
	Let $A\in M_{n\times n}(\R)$ be symmetric. Then all eigenvalues of $A$ are real,
	and there is an orthogonal matrix $Q$ (that is, $Q^{\transp}Q=I$) and a diagonal
	matrix $\Lambda=\diag(\lambda_1,\dots,\lambda_n)$ with
	$A=Q\Lambda Q^{\transp}$. Equivalently, $\R^n$ has an orthonormal basis of
	eigenvectors of $A$.
\end{theorem}

\begin{proof}[Proof sketch]
	Reality of the eigenvalues: if $Av=\lambda v$ with $v\in\C^n\setminus\{0\}$,
	then $\bar v^{\transp}Av=\lambda\bar v^{\transp}v$; taking conjugate transposes
	and using $A=A^{\transp}=\bar A$ gives
	$\bar v^{\transp}Av=\bar\lambda\bar v^{\transp}v$; since
	$\bar v^{\transp}v=\norm{v}^2>0$ we get $\lambda=\bar\lambda$.

	Diagonalisation proceeds by induction on $n$. Let $\lambda_1$ be an eigenvalue
	with unit eigenvector $q_1$; the subspace $q_1^{\perp}$ is $A$-invariant,
	because $\ip{q_1}{Aw}=\ip{Aq_1}{w}=\lambda_1\ip{q_1}{w}=0$ for
	$w\perp q_1$. Applying the induction hypothesis to the restriction of $A$ to
	$q_1^\perp$, which is symmetric with respect to the induced inner product, and
	adjoining $q_1$ produces the orthonormal eigenbasis. The full argument, in the
	form of the real spectral theorem, is \textcite[Ch.~7]{axler2024linear}.
\end{proof}

\begin{definition}[Quadratic form and definiteness]\label{def:definiteness}
	Let $A$ be symmetric. The \dfn{quadratic form} associated with $A$ is
	$q_A(x)=x^{\transp}Ax$. The matrix $A$ is
	\dfnas{positive definite}{positive definite matrix} ($A\succ0$) if $q_A(x)>0$
	for all $x\neq0$; \dfnas{positive semidefinite}{positive semidefinite matrix}
	($A\succeq0$) if $q_A(x)\geq0$ for all $x$; and negative (semi)definite if $-A$
	is positive (semi)definite. Otherwise $A$ is \dfn{indefinite}.
\end{definition}

\begin{proposition}[Definiteness via the spectrum]\label{prop:definiteness-eigen}
	A symmetric $A$ is positive definite if and only if all its eigenvalues are
	strictly positive, and positive semidefinite if and only if all are
	non-negative. Moreover
	\[
		\lambda_{\min}\norm{x}^2\leq x^{\transp}Ax\leq\lambda_{\max}\norm{x}^2
		\qquad\text{for all }x\in\R^n,
	\]
	where $\lambda_{\min},\lambda_{\max}$ are the smallest and largest eigenvalues.
\end{proposition}

\begin{proof}
	Write $A=Q\Lambda Q^{\transp}$ as in \autoref{thm:spectral} and set
	$y=Q^{\transp}x$. Since $Q$ is orthogonal, $\norm{y}=\norm{x}$ and
	$x\neq0\iff y\neq0$, and
	\[
		x^{\transp}Ax=y^{\transp}\Lambda y=\sum_{i=1}^n\lambda_iy_i^2 .
	\]
	If every $\lambda_i>0$ this is strictly positive for $y\neq0$. Conversely, if
	$\lambda_j\leq0$ for some $j$, taking $y=e^j$, that is $x=Qe^j$, gives
	$x^{\transp}Ax=\lambda_j\leq0$ with $x\neq0$. The semidefinite case is
	identical with weak inequalities. The two-sided bound follows from
	$\lambda_{\min}\sum_iy_i^2\leq\sum_i\lambda_iy_i^2\leq\lambda_{\max}\sum_iy_i^2$.
\end{proof}

\begin{remark}[Only the symmetric part matters]\label{rem:symmetric-part}
	For any square $A$, $x^{\transp}Ax=x^{\transp}\bigl(\tfrac12(A+A^{\transp})
		\bigr)x$, because $x^{\transp}Ax$ is a scalar and hence equals its own
	transpose $x^{\transp}A^{\transp}x$. A quadratic form therefore determines and
	is determined by the symmetric part of $A$ alone, and one may always assume $A$
	symmetric. This is the content of the gradient computation in
	\autocite[slide 49]{fukuda2026calculus}: for $f(x)=x^{\transp}Ax$,
	\[
		\nabla f(x)=(A+A^{\transp})x,
	\]
	which reduces to $2Ax$ exactly when $A$ is symmetric. Forgetting the factor
	$A+A^{\transp}$ is one of the most common errors in matrix differentiation.
\end{remark}

\begin{proposition}[Leading principal minors]\label{prop:sylvester}
	Let $A$ be symmetric $n\times n$ and let $D_k$ denote the determinant of its
	upper-left $k\times k$ submatrix. Then $A\succ0$ if and only if $D_k>0$ for
	$k=1,\dots,n$; and $A\prec0$ if and only if $(-1)^kD_k>0$ for $k=1,\dots,n$,
	that is, the leading principal minors alternate in sign starting negative.
\end{proposition}

\begin{remark}[The semidefinite case is different]\label{rem:semidefinite-minors}
	The analogous criterion for \emph{semi}definiteness requires all principal
	minors, not merely the leading ones, to be non-negative. The matrix
	$A=\begin{psmallmatrix}0&0\\0&-1\end{psmallmatrix}$ has $D_1=D_2=0$, so all
	leading principal minors are non-negative, yet $A$ is negative semidefinite and
	not positive semidefinite. This trap matters when checking second-order
	conditions at a boundary optimum in \autoref{ch:optimization}, where
	semidefiniteness rather than definiteness is what the theory delivers.
\end{remark}

\section{Economic Application: Linear Dynamics and Stability}\label{sec:lin-dynamics}

Consider the linear difference system
\begin{equation}\label{eq:linear-system}
	x_{t+1}=Ax_t,\qquad x_0\in\R^n\text{ given},
\end{equation}
which arises as the linearisation of a nonlinear dynamic model around a steady
state --- for instance the Ramsey model of \autoref{sec:dyn-ramsey} linearised
about $(k^\ast,c^\ast)$. Iterating, $x_t=A^tx_0$.

\begin{proposition}[Stability]\label{prop:linear-stability}
	For system~\eqref{eq:linear-system}: $x_t\to0$ for every $x_0\in\R^n$ if and
	only if $\rho(A)<1$. If $\rho(A)>1$, then $x_t$ diverges for every $x_0$ outside
	a proper subspace of $\R^n$.
\end{proposition}

\begin{proof}[Proof for the diagonalisable case]
	Suppose $A=P\Lambda P^{-1}$ with $\Lambda=\diag(\lambda_1,\dots,\lambda_n)$.
	Then $A^t=P\Lambda^tP^{-1}$ with
	$\Lambda^t=\diag(\lambda_1^t,\dots,\lambda_n^t)$. Writing
	$z_t=P^{-1}x_t$ decouples the system into $z_{t,i}=\lambda_i^tz_{0,i}$. Each
	coordinate tends to $0$ for every initial condition if and only if
	$\abs{\lambda_i}<1$, and grows without bound whenever $\abs{\lambda_i}>1$ and
	$z_{0,i}\neq0$. The set of $x_0$ with $z_{0,i}=0$ for all $i$ with
	$\abs{\lambda_i}>1$ --- the stable subspace --- is a proper subspace when some
	$\abs{\lambda_i}>1$. The non-diagonalisable case is handled by the Jordan form,
	where a $\lambda$-block of size $k$ contributes terms $t^{j}\lambda^{t}$ with
	$j<k$; these still tend to $0$ when $\abs{\lambda}<1$, so the criterion is
	unchanged.
\end{proof}

\begin{remark}[Saddle-path stability]\label{rem:saddle-path}
	In macroeconomic models the interesting case is neither $\rho(A)<1$ nor
	$\rho(A)>1$ but a mixture: some eigenvalues inside the unit circle and some
	outside. The state vector then splits into predetermined variables, whose
	initial values are given, and jump variables, whose initial values are chosen.
	A unique bounded equilibrium path exists when the number of eigenvalues outside
	the unit circle equals the number of jump variables --- the Blanchard--Kahn
	count. Economically, the jump variables must be placed on the stable subspace,
	and the counting condition says that the placement problem has exactly one
	solution, which by \autoref{cor:solvability}(iii) is a statement about a square
	invertible system. The recursive treatment is
	\textcite[Ch.~2]{ljungqvist2018recursive}.
\end{remark}

\begin{eg}[Leontief input--output]\label{eg:leontief}
	Let $A\in M_{n\times n}(\R_+)$ have $a_{ij}$ equal to the units of good $i$
	required per unit of good $j$, and let $d\in\R^n_+$ be final demand. Gross
	output $x$ must satisfy $x=Ax+d$, that is $(I-A)x=d$. A non-negative solution
	exists for every $d\in\R^n_+$ if and only if $(I-A)^{-1}$ exists and is
	non-negative, which holds if and only if $\rho(A)<1$; in that case
	\[
		(I-A)^{-1}=\sum_{k=0}^{\infty}A^k,
	\]
	the Neumann series, whose $k$-th term is the input requirement generated at the
	$k$-th round of production. The condition $\rho(A)<1$ is thus the productivity
	condition that the economy generates a surplus. Convergence of the series is a
	fixed-point statement: $x\mapsto Ax+d$ is a contraction in a suitable norm, and
	\autoref{thm:banach-fixed-point} applies. The same series reappears in
	\autoref{ch:dynamic} as the present-value expansion of a value function.
\end{eg}

\section{Chapter Summary}\label{sec:lin-summary}

A linear map between finite-dimensional spaces is completely described, once
bases are fixed, by a matrix whose columns are the images of the basis vectors
(\autoref{thm:matrix-rep}). Rank--nullity converts questions about solvability
into questions about dimension (\autoref{thm:rank-nullity},
\autoref{cor:solvability}); for square matrices, and only for square matrices,
injectivity, surjectivity, invertibility, and non-vanishing determinant coincide.
The determinant measures oriented volume magnification, which is why it appears
in the change-of-variables formula. Symmetric matrices are orthogonally
diagonalisable with real eigenvalues (\autoref{thm:spectral}), so definiteness of
a quadratic form is a statement about the sign of the spectrum
(\autoref{prop:definiteness-eigen}) --- the fact used for second-order conditions
throughout \autoref{ch:optimization}. Only the symmetric part of a matrix affects
its quadratic form, and $\nabla(x^{\transp}Ax)=(A+A^{\transp})x$. Stability of a
linear dynamic system is governed by the spectral radius.

\begin{exercise}\label{ex:lin-1}
	Verify directly, without invoking \autoref{rem:symmetric-part}, that for
	$f(x_1,x_2)=a_{11}x_1^2+(a_{12}+a_{21})x_1x_2+a_{22}x_2^2$ one has
	$\nabla f(x)=(A+A^{\transp})x$, and conclude that two distinct matrices can
	induce the same quadratic form. (This is Exercise~3.3 of
	\textcite{fukuda2026notes}; see \autocite[slide 49]{fukuda2026calculus}.)
\end{exercise}

\begin{exercise}\label{ex:lin-2}
	Let $X$ be $n\times k$ with $\rank X=k\leq n$. Show that $X^{\transp}X$ is
	symmetric positive definite, hence invertible. Where is the rank hypothesis
	used, and what fails without it? (The result is what makes the OLS formula
	$\hat\beta=(X^{\transp}X)^{-1}X^{\transp}y$ of
	\autocite[slide 60]{fukuda2026calculus} well posed.)
\end{exercise}

\begin{exercise}\label{ex:lin-3}
	Let $A\in M_{n\times n}(\R_+)$ with all row sums strictly less than $1$. Show
	$\rho(A)<1$, and deduce from \autoref{eg:leontief} that the Leontief system is
	productive. Interpret the row-sum condition economically.
\end{exercise}
